% =====================================================================
% Section 7 — Discussion and Threats to Validity
% =====================================================================
\section{Discussion and Threats to Validity}\label{sec:threats}

\subsection{Threats to validity}

\paragraph{Measurement environment.} All frameworks were measured under
identical container limits on a single bare-metal host, with the load
generator pinned to cores disjoint from the server under test, the CPU
governor fixed to \texttt{performance}, and Turbo Boost disabled to
remove frequency variance. This isolates the comparison from
client--server contention and thermal drift, but the results remain
\emph{relative}: absolute numbers would differ on other hardware, and a
containerised single-host set-up does not model network latency or
multi-node deployment. We mitigate run-to-run noise with five
repetitions, a warm-up phase that lets the JIT reach steady state, and
by reporting medians with observed ranges rather than single points.

\paragraph{Endpoint parity.} The two endpoints (\texttt{/plaintext} and
\texttt{/json}) are deliberately minimal and functionally equivalent
across frameworks, which isolates framework overhead but does not
exercise persistence, templating, or business logic. The throughput
figures should therefore be read as a measure of baseline request
handling, not of application-level performance.

\paragraph{Artifact-size comparison.} JxMVC's \SI{253}{\kilo\byte} is
the framework JAR with the servlet container excluded (it is a
\texttt{provided} dependency, \(\approx\)\SI{15}{\mega\byte} for
Tomcat/Jakarta~EE), whereas an über-JAR such as Spring Boot's bundles
its embedded server. We therefore report deployable \emph{image} size
in the evaluation for a like-for-like ``framework~+~server''
comparison, and are explicit that the \SI{253}{\kilo\byte} figure
refers to the auditable framework core, not to a self-contained
deployable.

\subsection{Limitations}

We state the design's limitations plainly. The most visible is
\textbf{resident memory}: JxMVC's RSS is the highest of the JVM stacks,
a direct consequence of running on an embedded Tomcat servlet
container; reducing it is future work. The framework uses
\textbf{reflection} in request binding and JSON handling, which trades
some hot-path cost and native-image friendliness for API simplicity
(partially mitigated by caching annotated-method metadata per class).
Convenience features carry caveats: \texttt{@JxQuery} derives context
from the call stack; SQL is composed by concatenation with
parameterised mitigations rather than a full query DSL; classpath
scanning must handle both \texttt{file:} and \texttt{jar:} layouts; the
scheduler uses five-field cron; and the event bus is synchronous. These
are conscious simplicity/functionality trade-offs, not oversights.
Finally, JxMVC is a young project with a single principal architect and
limited adoption relative to the incumbents; maturity and community are
themselves risks for production use.

\subsection{Why another framework?}

The natural objection to any new framework is that the space is
crowded. Our answer is neither performance supremacy nor feature
parity---the evaluation shows JxMVC does not lead on throughput and
carries a memory cost. Rather, JxMVC occupies an under-served
\emph{niche}: a full-stack, secure-by-default MVC framework small
enough to audit in full and to teach from, with zero runtime
dependencies. Its value is the combination of that design point, the
empirical evidence that it remains competitive on latency and start-up,
and its use as a readable reference implementation of web-framework
internals.

\subsection{Practical implications: when to choose JxMVC}

The evaluation and our experience suggest a concrete decision guide
rather than a blanket recommendation. JxMVC is a strong fit when the
\emph{supply-chain surface} must be minimised or fully audited---the
entire runtime is one \SI{253}{\kilo\byte} JAR with no transitive
dependencies to vet or patch, a property that matters in regulated,
air-gapped, or security-reviewed environments where each added library
is a compliance cost. It fits \emph{internal tools and services} where
auditability and safe defaults outweigh raw feature breadth, and
\emph{start-up-sensitive} deployments (scale-to-zero, frequent
redeploys) where its sub-second cold start beats Spring Boot and
Micronaut. Conversely, JxMVC is \emph{not} the right choice when
per-instance memory density is the binding cost (the servlet container
makes its RSS the highest of the field), when peak throughput is
paramount (Quarkus and Javalin lead), or when a team depends on a large
plugin ecosystem and commercial support---areas where the incumbents
are, and should remain, the default.

\subsection{Pedagogical value}

A recurring, non-performance benefit emerged repeatedly while building
and documenting the framework: JxMVC is small and explicit enough to
serve as a \emph{readable reference implementation} of web-framework
internals. Because routing, dependency injection, the connection pool,
JSON parsing, the security pipeline, and an OIDC client are each a few
hundred lines of dependency-free Java on top of the JDK, a student can
read a mechanism end to end in one sitting rather than tracing it
through layers of generated code and third-party abstractions. The
same property aids security review: the complete attack-relevant
surface---sanitisation, authentication ordering, token verification,
body caps---can be inspected directly. We therefore see JxMVC's
compactness not only as an operational advantage but as a teaching and
auditing artifact, turning the web stack from a black box into a worked
example.
